\documentclass{beamer}

\usepackage[utf8]{inputenc}
\usepackage[T1]{fontenc}
\usepackage{eurosym}
\usepackage{rotating}
\usepackage{comment}

\AtBeginSection[]
{
  \begin{frame}
  \frametitle{Wir sind hier:}
    \tableofcontents[currentsection]
  \end{frame}
}

\begin{document}

\title{Erkennen eines Graphen aus einer Punktwolke}
\date{}

\frame{\titlepage}

\section{Einleitung}
\frame{\frametitle{Das Problem}
\begin{description}
\item[Eingabe] ein metrischer Raum $(X, d)$, z.B.
  \begin{enumerate}
  \item GPS-Koordinaten eines Autos, das in einer Stadt herumfährt
  \item Koordinaten der schwarzen Pixel in einem Schwarz-Weiß-Bild
  \end{enumerate}
\item[Ausgabe] eine Approximation des metrischen Graphen, der dem metrischen Raum zugrunde liegt, z.B.
  \begin{enumerate}
  \item das Straßennetz der Stadt
  \item der Graph, der im Bild abgebildet wird\\ $\rightarrow$ Schrifterkennung
  \end{enumerate}
\end{description}
}

\frame{\frametitle{Wozu ist das nützlich?}

\begin{enumerate}
\item Struktur in gro{\ss}e Mengen geometrischer Daten bringen um deren Analyse zu ermöglichen/erleichtern \pause
\item Häufig enthalten die Daten rauschen und sind umfangreich $\rightarrow$ Ziel: Datenmenge kompakt in ihren wichtigsten Verzweigungen darstellen
\end{enumerate}
}

\frame{\frametitle{Beispiel (Erdbeben)}
\begin{figure}
\includegraphics{images/faultlines.png}
\end{figure}
}

\frame{\frametitle{Beispiel (Straßennetz)}
\begin{figure}
\includegraphics{images/roadnetwork.png}
\end{figure}
}

\section{Vorverarbeitung}

\frame{\frametitle{Vorverarbeitung}
\begin{enumerate}
\item Menge von Punkten erstellen
  \begin{itemize}
  \item Koordinaten aus GPS-Spur-Dateien extrahieren
  \item Koordinaten der schwarzen Pixel eines Bildes bestimmen
  \end{itemize}
\item ggf. Punktmenge auf repräsentative Teilmenge reduzieren
  \begin{itemize}
  \item metrisches $\varepsilon$-Netz
  \item Laufzeit: $O(n)$ (unter Verwendung eines Gitters)
  \end{itemize}
\item auf die Punktmenge basierenden $\alpha$-Komplex erstellen
  \begin{itemize}
  \item $\alpha$-Komplex: Unterkomplex der Delaunay-Triangulierung
  \item lange Kanten der Triangulierung werden entfernt
  \item verwenden JTS Topology Suite (DelaunayTriangulationBuilder)
  \end{itemize}
\item Abstandsmethode implementieren
  \begin{itemize}
  \item $d(x,y) :=$ Länge des kürzesten Pfades zwischen $x$ und $y$
  \item Algorithmus von Floyd-Warshall
  \end{itemize}
\end{enumerate}
}

\frame{\frametitle{Klassendiagramm}
\begin{figure}
\includegraphics[height=8cm]{images/preprocessing.png}
\end{figure}
}

\frame{\frametitle{Verbesserungsmöglichkeiten}
\begin{itemize}
\item $\alpha$-Komplex
  \begin{itemize}
  \item um Lücken zu vermeiden, muss man u.\,U. ein großes $\alpha$ wählen
  \item führt dazu, dass mehr Knoten durch Kanten verbunden werden\\ $\rightarrow$ Struktur schwieriger zu erkennen, Abstände verzerrt
  \item Idee: Zeitstempel in den GPS-Spur-Dateien in Betracht ziehen
  \end{itemize}
\item Laufzeit der Abstandsberechnung
  \begin{itemize}
  \item Floyd-Warshall Laufzeit: $O(n^3)$
  \item $\alpha$-Komplexe enthalten $O(n)$ Kanten
  \item Idee: Algorithmus von Dijkstra verwenden $\rightarrow ~ O(n^2\log n)$
  \end{itemize}
\item Abstandsberechnung für GPS-Koordinaten
  \begin{itemize}
  \item aktuell interpretieren wir die Koordinaten als Punkte auf einer Ebene und berechnen den euklidischen Abstand
  \item Idee: Orthodrome berechnen $\rightarrow$ Abstand in km
  \end{itemize}
\end{itemize}
}

\section{Der Algorithmus}

\frame{\frametitle{Drei wesentliche Schritte}
\begin{enumerate}
	\item[1] \textbf{Labeling}: Welche Punkte gehören zu Kanten, welche zu Knoten? $\rightarrow$ Punkte mit entsprechendem Label versehen\pause
	\item[2] \textbf{Rekonstruktion}: Anhand der Labels bestimmen, welche Punkte sich zu einem Knoten und welche sich zu einer Kante zusammentun $\rightarrow$ Rekonstruktion des der Punktmenge zugrundeliegenden Graphen\pause
	\item[3] \textbf{Metrik wiederherstellen}: Kanten mit Abständen versehen
\end{enumerate}
}

\frame{\frametitle{Zu Schritt 1: Labeling}
\begin{enumerate}
\item Betrachten einen kreisförmigen Ausschnitt um jeden Punkt und bestimmen Anzahl der Zusammenhangskomponenten darin
\begin{figure}
\includegraphics[height=2cm,width=2cm,keepaspectratio]{images/branch.png}
\end{figure}
	\begin{itemize}
		\item Grad des Punkts = \# Zusammenhangskomponenten im kreisförmigen Ausschnitt
	\end{itemize}\pause
\item Grad == 2 $\rightarrow$ \textit{edge point} (wird später zu einer Kante gehören)\pause
\item Grad != 2 $\rightarrow$ \textit{preliminary branch point} (wird später zu einem Knoten gehören, aber noch umbenannt zu \textit{branch point})\pause
\item Alle weniger als $x$ von einem \textit{preliminary branch point} entfernten Punkte werden als \textit{branch points} eingeordnet (Knoten)
\end{enumerate}
}

\frame{\frametitle{Zu Schritt 2: Rekonstruktion des Graphen}
\begin{itemize}
	\item Mithilfe der Labels bestimmen, welche Punkte sich zu einem Knoten und welche sich zu einer Kante zusammentun\pause
	\begin{itemize}
		\item Punkte sind jetzt nach dem Labeling in zwei Mengen aufgeteilt:
		\begin{itemize}
		\item \textit{edge points}
		\item \textit{branch points}
		\end{itemize}\pause
\item Um die Anzahl der späteren Kanten und Knoten zu bestimmen, wird Rips-Vietoris-Graph für beide Mengen erstellt\pause
	\begin{itemize}
		\item Im Rips-Vietoris-Graphen werden alle Punkte, die innerhalb eines gewissen Abstands zueinander liegen, zu einer Zusammenhangskomponente gefasst
    \end{itemize}\pause
\item Jede Zusammenhangskomponente entspricht einem Knoten, bzw. einer Kante\pause
\item Zwei Punkte werden durch eine Kante verbunden, wenn sie Punkte in ihrer Zusammenhangskomponente haben, die einen gewissen Abstand zu Punkten in der selben Zusammenhangskomponente einer Kante nicht überschreiten
	\end{itemize}
\end{itemize}
}

\frame{\frametitle{Zu Schritt 3: Metrik wiederherstellen}
\begin{itemize}
	\item Kanten mit Abständen versehen
	\begin{itemize}
		\item Jeder Kante wird als Länge der Durchmesser (längster kürzester Weg) ihrer Zusammenhangskomponente zugewiesen
	\end{itemize}
\end{itemize}
}

\frame{\frametitle{Klassendiagramm}
\begin{figure}
\includegraphics[width=11cm]{images/main.png}
\end{figure}
}

\section{Organisation}

\frame{\frametitle{Kommunikation und Tools}
\begin{itemize}
\item Kommunikation und Tools
	\begin{itemize}
	\item regelmäßige Treffen
	\item Skype
	\item E-Mail
	\item Trello
	\item GitHub
	\end{itemize}
\item Programmierung
	\begin{itemize}
	\item Festlegung von Interfaces
	\item 3 Gruppen
	\end{itemize}
\end{itemize}
}

\section{Probleme und Lernerfolge}

\frame{\frametitle{Probleme und Lernerfolge}
\begin{itemize}
\item Probleme
	\begin{itemize}
	\item Testen der Komponenten schwierig
	\item Absprung eines Gruppenpartners
	\end{itemize}
\item Lernerfolge
	\begin{itemize}
	\item Arbeit in Teilgruppen
	\item Planung der Software
	\item Zusammenführen der Teilstücke
	\end{itemize}
\end{itemize}
}

\end{document}
